\par\bigskip
\hypertarget{native-379-parabolic-remarks}{}\label{native-379-parabolic-remarks}
\pdfbookmark[2]{Bemerkungen zur parabolischen Bewegung}{parabolic-bemerkungen}
\begin{center}\textsc{Bemerkungen.}\end{center}
\begingroup\small
\noindent\hspace*{\parindent}In der »Zeitschrift für Astronomie, herausgegeben von \textsc{Lindenau} und \textsc{Bohnenberger}«, Erster Band, steht in der von \textsc{Lindenau} im November 1815 entworfenen Einleitung, S. 46, die Bemerkung: »\textsc{Gauss} hat neuerlich beinahe die ganze Cometentheorie umgearbeitet, neue Tafeln entworfen, und wir haben ein eigenthümliches Werk darüber von ihm zu erwarten, was wohl nichts zu wünschen übrig lassen wird.« Die Schlussworte bilden eine von \textsc{Gauss} häufig gebrauchte Redewendung, dürften also einem Briefe von \textsc{Gauss} an \textsc{Lindenau} entnommen sein. Hiemit in Übereinstimmung steht der Brief an \textsc{Olbers} vom 13. Jan. 1815, S. 335, wo sich dieselbe Redewendung findet. In einem Brief an \textsc{Schumacher} vom 21. Nov. 1825 sagt \textsc{Gauss}: »Der Wunsch, den ich immer bei meinen Arbeiten gehabt habe, ihnen eine solche Vollendung zu geben, ut nihil amplius desiderari possit, erschwert sie mir freilich ausserordentlich.« Man sehe auch \textsc{Schumachers} Antwort darauf vom 2. Dec. 1825. Dass \textsc{Gauss} sich namentlich im December 1814 und im Januar 1815 mit der Bahnbestimmung der Cometen beschäftigt hat, geht auch aus dem Briefwechsel mit \textsc{Olbers} (S. 332 u. 350) hervor, während die in Band VI, Seite 25 ff. abgedruckte Abhandlung über diesen Gegenstand bereits im September 1813 der Göttinger Societät vorgelegt worden ist.
\par
\SourceLeaf{380}
\noindent\hspace*{\parindent}Auf Seite 325—370 sind alle Notizen zusammengestellt, welche sich im Nachlass darüber vorgefunden haben. Die Notizen [1] und [2] stehen in Handbüchern, [3] auf der Einbandseite eines Exemplars von »\textsc{Lambert}, insigniores orbitae cometarum proprietates, Augustae Vindelicorum 1761«; [4] ist aus Eintragungen im vorgenannten Buche und einem Handbuche zusammengestellt. Die Fussnote auf S. 336 stammt aus dem Umschlag eines Exemplars der deutschen Übersetzung der oben erwähnten Band VI, S. 25 ff. abgedruckten Abhandlung, welches auch viele andere hieher gehörige Notizen enthält (z. B. eine andere Ableitung der Formel \(\dfrac{\text{Ausschnitt}}{\text{Dreieck}}=\dfrac{2+\cos\varphi}{3\cos\varphi}\mathpunct{\text{\raisebox{.45ex}{,}}}\) als die S. 329—330 gegebene), deren Inhalt sich aber durch die abgedruckten Notizen und Briefstellen erledigt. Ausser der S. 331 abgedruckten Tafel finden sich Andeutungen über weitere Hülfstafeln, die \textsc{Gauss} zur Lösung der auftretenden Gleichungen zu entwerfen beabsichtigte.
\par
\noindent\hspace*{\parindent}Merkwürdigerweise ist von den in den vorgenannten Notizen benutzten Kunstgriffen in der Notiz [7], welche im Januar 1815 entstanden zu sein scheint und vermuthlich das im Brief an \textsc{Olbers} vom 13. Januar 1815 (S. 337) erwähnte »Musterbeispiel« enthält, nur weniges wiederzufinden. Diese Notiz steht in einem Handbuche, die Bezeichnungen sind wohl verständlich, wenn sie auch zum Theil von denen der Abhandlung vom September 1813 abweichen. Übrigens ist in dieser Notiz ein Rechenfehler untergelaufen: \textsc{Gauss} hatte (vgl. Seite 341) den Werth von \(Q\) zu \(272^0\ 40'\ 21''\) gefunden statt \(294^0\ 31'\ 57''\) und hat auch diesen Fehler selbst später bemerkt. Er hatte beabsichtigt, dass Beispiel mit dem Werthe \(P=Q\) weiter zu rechnen, weil diese Wahl von \(P\) die grösste Schärfe (vgl. Seite 341) bietet. Infolge des Rechenfehlers ist aber die weitere Rechnung mit \(P=272^0\ 40'\ 21''\) (S. 342) ausgeführt, und in dieser Form bei der Herausgabe beibehalten, weil ja doch die Wahl von \(P\) beliebig und die Abweichung von dem richtigen Werthe von \(Q\) nicht sehr erheblich ist.
\par
\noindent\hspace*{\parindent}Zu der Seite 351—367 abgedruckten Tafel zur parabolischen Bewegung vergleiche man die Fussnote in der Theoria motus S. 56 dieses Bandes. Gauss hat diese Tafel, wie er selbst (S. 353) sagt, auf 5 Decimalen der Bogensecunde berechnet und beabsichtigte, sie in Intervalle von \(\frac1{1000}\) für das Argument \(\log(At)\) und auf 2 Decimalen der Bogensecunde abzudrucken; aus diesem Grunde ist bei der Herausgabe die 3. Decimale beibehalten worden, so dass sich die Interpolation (vgl. S. 371) mit der von \textsc{Gauss} gewünschten Schärfe für das von ihm beabsichtigte Intervall ausführen lässt.
\par
\noindent\hspace*{\parindent}Die Tafel steht so, wie sie abgedruckt ist, in einem Handbuch, zum Theil auch auf einem Convolut loser Blätter; an ersterer Stelle trägt sie die Unterschrift »Vollendet Febr. 25. 1816«, an letzterer: »Vollendet 1816 Febr. 21.«
\par
\noindent\hspace*{\parindent}Das hinter der Tafel abgedruckte Schreiben, das sich im Nachlass vorfand, war offenbar an \textsc{Encke} gerichtet, der hier rechnen half; es verdient besonders deswegen Interesse, weil es die ganze Technik zeigt, mit der \textsc{Gauss} seine Tafeln zu berechnen pflegte.
\par
\noindent\hspace*{\parindent}Auffallend ist übrigens, dass der S. 371 abgedruckte im Jahre 1843 veröffentlichte Artikel gar keine Andeutung mehr von \textsc{Gauss} eigener Tafel enthält.
\par\smallskip
\hfill\textsc{Brendel}.
\par\endgroup\bigskip
\begin{center}\rule{.27\linewidth}{.4pt}\end{center}
